%==============================================================================
% 三、模型假设
%==============================================================================
\section{模型假设}

本文只保留五条真正影响建模取法的假设。每条给出数据层面的依据，并注明它被哪一节用到；凡依据可实测的都附实测值。其余常规假设（如随机误差零均值）不单独列出。

\textbf{假设 1（量纲与缺失可分辨）}\quad 附件1 中经对数变换的化学列底数为 10 且可逐列判定；取值为地板值时表示低于检出限而非真实浓度，NaN 表示未测量。依据：附件8 的 Calibrated 表给出 70 个化学列的还原浓度，逐列对拍 37 列落在 $\log_{10}$ 域、33 列落在原始线性域、161 列无对照；地板值 $-1.817819$ 在 142 列上精确重复共 6807 次，而 NaN 只分布在 9 列、共 73 格，两者形态不同。\emph{用于}第五章：不做整列反变换、不删列，地板值保留、NaN 按列中位数填补。

\textbf{假设 2（感官强度连续）}\quad 附件4 的 50 个感官项是 7 点量表按品评员 z-score 后再按啤酒平均的连续强度，项间差异不含量纲差异。依据：50 列全部存在负下限，每列不同取值数中位 229，逐列仿射搜索的残差 0.11--0.15，落不到整数格点；平移后直接取整会使 53.2\% 的样本压到 0 档。\emph{用于}第 6.2 节：问题一按回归而非分类建模，0--6 分级留到第 6.8 节用分位映射处理。

\textbf{假设 3（样本独立、风格组内同质）}\quad 每款啤酒的化学谱与感官评分来自同一批次且样本间独立；同一风格内的化学谱可由单一多元分布刻画。依据：beer 列在附件1、4、7 中均唯一且 250 行顺序一致；对 22 个风格做 Kruskal--Wallis 检验，231 个特征中 205 个在 0.05 水平显著、148 个过 Bonferroni 校正，$\varepsilon^2$ 中位 0.1768。风格内样本数中位 11、最小 3，故风格层面的结论不外推到单个品牌。\emph{用于}第 6.3 节的交叉验证划分与第 8.2 节的风格内条件分布。

\textbf{假设 4（局部线性、代理只定向不定点）}\quad 在实测样本覆盖范围内化学谱到感官的映射可用线性结构近似，但线性代理只支撑“往哪个方向调”，不支撑“哪个点更好”。依据：嵌套交叉验证中弹性网在 6 项取胜、两类树模型在 44 项取胜，说明线性结构存在但非全局；偏最小二乘代理的绝对误差中位 0.4626、95 分位 1.4124、上界 3.1525，与相邻 Pareto 点之间的间隙同量级。\emph{用于}第 7.3 节改用树模型加逐列置换做归因，第 8.3 节用偏最小二乘的线性映射写入规划，第 8.5 节对前沿只做分段报告。

\textbf{假设 5（阈值判据有效）}\quad 缺陷化合物是否恶化，以其浓度与香气阈值之比是否上升为准，与阈值的绝对量级无关。依据：231 个特征中 43 个在附件8 或公开文献中有阈值，同单位直接相除不做量纲换算；余下 188 个无阈值，标为 N 档且不作化学结论。\emph{用于}第 7.4 节的气味活性值裁判（A/B/C/N 四档）与第 8.4 节缺陷物的软约束与上限约束。
